\subsection{整闭整环的例子}

命题 10. 每个主理想整环都是整闭的。\par

设 A 是主理想整环，K 是它的分式域，x 是 K 的元素。存在 A 的两个互素元素 a, b 使得 \(x = ab^{-1}\) （《代数》，第 VII 章，§1，第 2 节，命题 1 及第 VI 章，§1，第 11 节，命题 9 (DIV)）。若 x 在 A 上是整的，则它是一个多项式的根

\(X^{n} + c_{1}X^{n-1} + \cdots + c_{n}\) （A{[}X{]} 中的）。于是 \(a^{n} = b(-c_{1}a^{n-1} - \cdots - c_{n}b^{n-1})\) ，这证明 \(b\) 整除 \(a^{n}\) 。由于 \(a\) 与 \(b\) 互素，这蕴涵 \(b\) 在 A 中可逆（《代数》，第 VI 章，§ 1，第 12 节，命题 11 的推论 1 (DIV)）；从而 \(x \in A\) 。

引理 2. 设 R 是环，P 是 R{[}X{]} 中的首一多项式。存在一个含 R 的环 R'，使得在多项式环 R'{[}X{]} 中多项式 P 是次数为 1 的首一多项式的乘积。

我们对 \(n\) （\(\mathbf{P}\) 的次数）作归纳，引理对 \(n = 0\) 与 \(n = 1\) 是显然的。于是设 \(n > 1\) 。设 \(\mathfrak{a}\) 是 \(\mathbb{R}[\mathbf{X}]\) 中由 \(\mathbf{P}\) 生成的理想，设 \(f\) 是 \(\mathbb{R}[\mathbf{X}]\) 到 \(\mathbf{B} = \mathbb{R}[\mathbf{X}] / \mathfrak{a}\) 上的典范同态。由于 \(\mathbf{P}\) 是首一的，对每个多项式 \(\mathbf{Q} \in \mathbb{R}[\mathbf{X}]\) ，\(\deg(\mathrm{PQ}) = \deg(\mathrm{P}) + \deg(\mathrm{Q})\) ，由此 \(\mathfrak{a} \cap \mathbb{R} = 0\) ；\(f\) 在 \(\mathbb{R}\) 上的限制因而是单射。把 \(\mathbb{R}\) 与子环 \(f(\mathbb{R})\) （\(\mathbb{B}\) 的）借助 \(f\) 等同起来，并写 \(b = f(\mathbf{X})\) ，我们看到 \(b\) 是 \(\mathbf{P}\) 在 \(\mathbb{B}\) 中的根，这里 \(\mathbf{P}\) 视为 \(\mathbb{B}[\mathbf{X}]\) 中的多项式。于是存在首一多项式 \(\mathbf{Q}\) （在 \(\mathbb{B}[\mathbf{X}]\) 中，次数为 \(n - 1\)）使得 \(\mathrm{P}(\mathbf{X}) = (\mathbf{X} - b)\mathrm{Q}(\mathbf{X})\) （《代数》，第 IV 章，§ 1，第 4 节，命题 5）。由归纳假设，存在环 \(\mathbb{R}' \supseteq \mathbb{B}\) 使得在 \(\mathbb{R}'[\mathbf{X}]\) 中多项式 \(\mathbf{Q}\) 是次数为 1 的首一多项式的乘积；显然在 \(\mathbb{R}'[\mathbf{X}]\) 中 \(\mathbf{P}\) 这时是次数为 1 的首一多项式的乘积。

命题 11. 设 A 是环，R 是 A-代数，P 与 Q 是 R{[}X{]} 中的首一多项式。若 PQ 的系数都在 A 上是整的，则 P 与 Q 的系数都在 A 上是整的。

两次应用引理 2，我们看出存在一个环 \(\mathbf{R}'\) （含 \(\mathbf{R}\)）以及元素族 \((a_i)_{1 \leqslant i \leqslant m}, (b_j)_{1 \leqslant j \leqslant n}\) （\(\mathbf{R}'\) 的），使得在 \(\mathbf{R}'[\mathbf{X}]\) 中 \(\mathrm{P}(\mathbf{X}) = \prod_{i=1}^{m} (\mathbf{X} - a_i)\) ，\(\mathbf{Q}(\mathbf{X}) = \prod_{j=1}^{n} (\mathbf{X} - b_j)\) ；\(\mathbf{PQ}\) 的系数属于整闭包 \(\mathbf{A}'\) （\(\mathbf{A}\) 在 \(\mathbf{R}'\) 中的），因而（第 2 节，命题 7）元素 \(a_i\) （ \(1 \leqslant i \leqslant m\) ）与 \(b_j\) （ \(1 \leqslant j \leqslant n\) ）属于 \(\mathbf{A}'\) 。由此得出 \(\mathbf{P}\) 与 \(\mathbf{Q}\) 的系数在 \(\mathbf{A}\) 上是整的（第 1 节，命题 4 的推论 1）。

设 A 是整环，K 是它的分式域，\(K'\) 是 K-代数（未必交换）。给定一个在 K 上为代数的元素 \(x \in K'\) ，多项式 \(P \in K[X]\) （满足 \(P(x) = 0\) 者）组成一个理想 \(a \neq 0\) （\(K[X]\) 的） ，它必然是主理想（《代数》，第 IV 章，§1，第 5 节，命题 7）。存在唯一的首一多项式生成 a；把《代数》第 V 章，§3，第 1 节，定义 3 中引入的术语加以推广，这个首一多项式将称为 x 在 K 上的极小多项式。

推论. 设 A 是整环，K 是它的分式域，x 是 K-代数 \(K'\) （未必交换）的元素。若 x 在 A 上是整的，则 x 在 K 上的极小多项式 P 的系数在 A 上是整的（从而若 A 是整闭的，它们属于 A）。

由假设（第 1 节，定理 1）存在一个首一多项式 \(Q \in A[X]\) 使得 \(Q(x) = 0\) 。由于 P 在 K{[}X{]} 中整除 Q，由命题 11，P 的系数在 A 上是整的。

设 A 是环，R 是（交换）A-代数；定义 R 上 A-代数结构的同态 \(\phi: A \rightarrow R\) 可以唯一地扩张为 A 与 R 上的多项式环的同态 \(A[X] \rightarrow R[X]\) ，它保持 X 不变，从而 \(R[X]\) 具有一个典范的 A{[}X{]}-代数结构。

命题 12. 设 A 是环，R 是 A-代数，P 是 \(R[X_{1},\ldots,X_{n}]\) 中的多项式。要使 P 在 \(A[X_{1},\ldots,X_{n}]\) 上是整的，必须且只需 P 的系数在 A 上是整的。

把 \(R[X_{1},\ldots,X_{n}]\) 中的多项式视为关于 \(X_{n}\) 的多项式（系数在 \(R[X_{1},\ldots,X_{n-1}]\) 中），我们立即看出可化归为对 n=1 证明本命题。设 P 是 \(R[X]\) 中的多项式；由第 1 节命题 5 立即得出，若 P 的系数属于 A 在 R 中的整闭包 B，则属于 \(B[X]=B\otimes_{A}A[X]\) 的元素 P 在 A{[}X{]} 上是整的。反之，设 P 在 A{[}X{]} 上是整的，并设

\[
\mathrm{Q} (\mathrm{Y}) = \mathrm{Y} ^ {m} + \mathrm{F} _ {1} \mathrm{Y} ^ {m - 1} + \dots + \mathrm{F} _ {m}
\]

是系数为 \(F_{i} \in A[X]\) 且以 P 为根的首一多项式。设 r 是一个严格大于多项式 P 与 \(F_{i}\) （ \(1 \leqslant i \leqslant m\) ）的所有次数的整数，我们写

\[
\mathrm{P} _ {1} (\mathrm{X}) = \mathrm{P} (\mathrm{X}) - \mathrm{X} ^ {r}.
\]

于是 \(\mathbf{P}_1\) 是多项式

\[
\mathrm{Q} _ {1} (\mathrm{Y}) = \mathrm{Q} (\mathrm{Y} + \mathrm{X} ^ {r}) = \mathrm{Y} ^ {m} + \mathrm{G} _ {1} \mathrm{Y} ^ {m - 1} + \dots + \mathrm{G} _ {m}
\]

系数在 A{[}X{]} 中；因此我们可以写

\[
- \mathrm{P} _ {1} (\mathrm{P} _ {1} ^ {m - 1} + \mathrm{G} _ {1} \mathrm{P} _ {1} ^ {m - 2} + \dots + \mathrm{G} _ {m - 1}) = \mathrm{G} _ {m}.\tag{1}
\]

现在 \(r\) 的选取蕴涵 \(-\mathbf{P}_1\) 是 \(\mathbb{R}[\mathbf{X}]\) 中的首一多项式，\(\mathbf{G}_m(\mathbf{X}) = \mathbf{Q}(\mathbf{X}^r)\) 亦然，因为多项式 \(\mathrm{F}_k(\mathbf{X})\mathbf{X}^{r(m - k)}\) 的次数都 \(< rm\) （对 \(k \geqslant 1\) ）。我们首先断定多项式

\[
\mathrm{P} _ {1} ^ {m - 1} + \mathrm{G} _ {1} \mathrm{P} _ {1} ^ {m - 2} + \dots + \mathrm{G} _ {m - 1}
\]

（R{[}X{]} 中的）也是首一的；此外，由于 \(G_{m}\) 的系数属于 A，命题 11 表明 \(P_{1}\) 的系数在 A 上是整的，因而 P 的系数当然在 A 上是整的。

命题 13. 设 A 是环，R 是 A-代数，A' 是 A 在 R 中的整闭包。则 A{[}X₁, \ldots, Xₙ{]} 在 R{[}X₁, \ldots, Xₙ{]} 中的整闭包等于 A'{[}X₁, \ldots, Xₙ{]}。

这由命题 12 及第 2 节的定义 3 得出。

推论 1. 设 A 是整环，A' 是它的整闭包。则多项式环 A{[}X₁, \ldots, Xₙ{]} 的整闭包是 A'{[}X₁, \ldots, Xₙ{]}。

对 n 作归纳论证，问题立即化归为 n = 1 的情形。设 K 是 A 的分式域，它也是 A' 的分式域；若 A{[}X{]} 的分式域 K(X) 的元素 P 在 A{[}X{]} 上是整的，则它属于多项式环 K{[}X{]} ，因为后者是主理想整环（《代数》，第 IV 章，§ 1，第 5 节，命题 7）从而是整闭的（命题 10）；于是推论由应用于 R = K 的命题 13 得出。

推论 2. 设 A 是整环。要使多项式环 A{[}X \(_{1}\) , \ldots, X \(_{n}\) {]} 是整闭的，必须且只需 A 是整闭的。

推论 3. 若 K 是域，则每个多项式代数 \(\mathbf{K}[\mathbf{X}_1,\ldots ,\mathbf{X}_n]\) 都是整闭的整环。

